% Generated draft: requires mathematical and source review.
\documentclass[11pt]{article}
\usepackage[margin=25mm]{geometry}
\usepackage{fontspec}
\setmainfont{Latin Modern Roman}
\usepackage{amsmath,amssymb,mathtools,unicode-math}
\setmathfont{Latin Modern Math}
\usepackage{enumitem,needspace,xurl,hyperref}
\setlength{\parindent}{0pt}
\setlength{\parskip}{5pt}
\setlength{\emergencystretch}{3em}
\newcommand{\N}{\mathbb{N}}
\newcommand{\Z}{\mathbb{Z}}
\newcommand{\Q}{\mathbb{Q}}
\newcommand{\R}{\mathbb{R}}
\newcommand{\C}{\mathbb{C}}
\newcommand{\F}{\mathbb{F}}
\newcommand{\A}{\mathbb{A}}
\newcommand{\PP}{\mathbb{P}}
\newcommand{\E}{\mathbb{E}}
\newcommand{\eps}{\varepsilon}
\newcommand{\dd}{\,\mathrm{d}}
\begin{document}
\section*{UnsolvedMath ID: 30003491}
\textbf{Source problem code:} catalogue entry
\subsection{Definitions and mathematical statement}
Write \(k\) for the unspecified base field intended by the source, and write \(G_a\) and \(G_m\) for the additive and multiplicative algebraic groups over that field. Thus, on a commutative \(k\)-algebra \(R\), their groups of points are \(G_a(R)=(R,+)\) and \(G_m(R)=R^\times\). This notation does not impose a universal quantifier over base fields.

The supplied background describes \(\mathcal C\) as the full category of commutative algebraic groups. It does not specify the base field, its characteristic, the conventions governing the objects of \(\mathcal C\), or the exact structure used to define extensions. These choices remain unresolved. Under the Yoneda interpretation, \(\operatorname{Ext}^i_{\mathcal C}(G_a,G_m)\), for \(i\geq 1\), is the group of equivalence classes of exact extensions
\[
0\longrightarrow G_m\longrightarrow E_{i-1}\longrightarrow\cdots\longrightarrow E_0\longrightarrow G_a\longrightarrow 0,
\]
with terms in the intended category and equivalence generated by morphisms of extensions inducing the identity on the endpoints; addition is the Baer sum. This interpretation is conditional on the intended category and extension theory supporting this construction.

With those source conventions fixed, the requested problems are:
\begin{enumerate}
\item Describe \(\operatorname{Ext}^2_{\mathcal C}(G_a,G_m)\) by explicit generators and defining relations, in the proposed analogue of the Milnor conjecture.
\item Determine whether there exists an integer \(N\geq 1\) such that
\[
\forall i\in\Z,\qquad i\geq N\ \Longrightarrow\ \operatorname{Ext}^i_{\mathcal C}(G_a,G_m)=0.
\]
\end{enumerate}
The proposed Milnor-conjecture analogue is not described in the supplied record, so its intended generators, relations, and comparison objects cannot be specified here. The scope of the base-field assumptions and any uniformity requirement on \(N\) are likewise unspecified.
\subsection{Short English statement}
For the intended category of commutative algebraic groups, give generators and relations for the second extension group from the additive group to the multiplicative group, and ask whether all sufficiently high extension groups between these objects vanish. The base field, precise category, extension theory, and proposed Milnor-conjecture analogue require clarification.
\subsection{Sources}
\begin{itemize}
\item Local UnsolvedMath registry: supplied record with permanent ID \(30003491\), title \emph{Higher Extensions Between Additive and Multiplicative Algebraic Groups}, source problem code \texttt{catalogue entry}, in the Oberwolfach Reports Open Problems collection. The record gives proposed year 2017 and names no proposer.
\item Supplied report citation: \emph{Algebraic Groups (2018)}, Oberwolfach Reports, pp.~1281--1347, DOI \texttt{10.4171/owr/2017/21}. \url{https://doi.org/10.4171/owr/2017/21}. The literature-source label is ``Oberwolfach Reports 2018.'' The supplied background describes these reports as records of mathematical workshops and their open-problem discussions.
\item Supplied literature source: Brion, ``isogeny categories.'' \url{https://doi.org/10.4171/DM/576}.
\item Supplied literature source: Brion, ``homological dimension.'' \url{https://arxiv.org/abs/1709.00175}.
\item Literature assessment dated exactly \texttt{2026-08-21}: the supplied curation labels the problem open and states that the higher-extension question for the full category of commutative algebraic groups remains open, while later work establishes hereditary behavior only for related isogeny categories. This is supplied dated curation derived from the OWR material, not an independently verified literature assessment; its time-sensitive status and cited claims require independent verification.
\item The registry supplies the statement-status label \texttt{corrected\_verified} and explains that mathematical content was preserved while notation, typography, and sentence structure were repaired for publication. This is a supplied editorial qualification, not independent mathematical verification of this draft.
\end{itemize}
\end{document}
